\documentclass[11pt,letterpaper,oneside]{report}
\usepackage[utf8]{inputenc}
\usepackage[T1]{fontenc}
\usepackage[spanish,es-nodecimaldot]{babel}
\usepackage{geometry}
\usepackage{xcolor}
\usepackage{graphicx}
\usepackage{booktabs}
\usepackage{longtable}
\usepackage{tabularx}
\usepackage{array}
\usepackage[hidelinks]{hyperref}
\usepackage{fancyhdr}
\usepackage{enumitem}

\geometry{margin=2.4cm}
\definecolor{amarillo}{HTML}{FFD426}
\definecolor{crema}{HTML}{FFF9EC}
\definecolor{verde}{HTML}{32D583}
\definecolor{negro}{HTML}{111111}
\definecolor{gris}{HTML}{F2F2F2}
\hypersetup{pdftitle={Compiscript MIPS Compiler IDE -- Proyecto 3},pdfauthor={DragonsSlayers 2.0}}
\pagestyle{fancy}\fancyhf{}\fancyhead[L]{Compiscript MIPS Compiler IDE}\fancyhead[R]{Proyecto 3}\fancyfoot[C]{\thepage}
\setlength{\headheight}{14pt}
\setlist{itemsep=2pt,topsep=4pt}
\newcommand{\codigo}[1]{\texttt{#1}}
\newcommand{\archivo}[1]{\texttt{\detokenize{#1}}}

\begin{document}

\begin{titlepage}
  \pagecolor{crema}\color{negro}
  \centering
  \vspace*{1.5cm}
  {\Large Universidad del Valle de Guatemala\par}
  \vspace{0.4cm}
  {\large Construcci\'on de Compiladores -- Secci\'on 10\par}
  \vspace{2cm}
  \fcolorbox{negro}{amarillo}{\parbox[c][2.4cm][c]{11cm}{\centering\bfseries\fontsize{25}{28}\selectfont COMPISCRIPT\\MIPS COMPILER IDE}}
  \vspace{1.4cm}

  {\Huge\bfseries Informe t\'ecnico\par}
  \vspace{0.3cm}
  {\LARGE Proyecto 3: generaci\'on de c\'odigo ensamblador\par}
  \vfill
  \begin{tabular}{rl}
    Equipo: & DragonsSlayers 2.0\\
    Integrantes: & Pablo Daniel Barillas Moreno\\
                 & Hugo Daniel Barillas Aj\'in\\
                 & Ernesto Ascencio\\
    Versi\'on: & 3.0.0\\
    Fecha: & 20 de octubre de 2026
  \end{tabular}
  \vfill
  {\large Compiscript $\rightarrow$ TAC $\rightarrow$ MIPS32 $\rightarrow$ MARS 4.5\par}
\end{titlepage}
\nopagecolor

\tableofcontents

\chapter{Resumen ejecutivo}

El Proyecto 3 completa un compilador educativo para Compiscript. Conserva el lexer y parser generados con ANTLR4TS, el an\'alisis sem\'antico con tabla de s\'imbolos y la representaci\'on intermedia de tres direcciones de las entregas anteriores. Sobre esa base incorpora un backend que transforma exclusivamente TAC en ensamblador MIPS32 compatible con MARS 4.5.

El producto se presenta como una interfaz gr\'afica React con editor Monaco. El usuario puede escribir o seleccionar un archivo \archivo{.cps}, ejecutar el pipeline, observar diagn\'osticos y examinar tokens, CST, \'arbol sem\'antico, s\'imbolos, \'ambitos, TAC, bloques b\'asicos, registros de activaci\'on y ensamblador. La exportaci\'on MIPS se habilita solamente cuando todas las fases terminan correctamente.

El backend implementa convenci\'on de llamadas mediante pila, asignaci\'on local de registros, aritm\'etica entera y flotante, control de flujo, funciones, llamadas anidadas, recursi\'on, arreglos, objetos, m\'etodos, strings y closures. La salida fue ensamblada y ejecutada en el MARS 4.5 oficial con quince programas positivos.

\chapter{Fundamentos te\'oricos}

\section{Fases de compilaci\'on}

Un compilador transforma un lenguaje fuente en una representaci\'on destino mediante fases con contratos definidos. En este proyecto, los tokens del lexer alimentan al parser; el CST v\'alido alimenta al visitor sem\'antico; el resultado tipado habilita TAC; finalmente, TAC habilita MIPS.

\begin{figure}[ht]
\centering
\small
\setlength{\fboxsep}{5pt}
\colorbox{amarillo}{\textbf{Compiscript}} $\longrightarrow$
\colorbox{amarillo}{\textbf{Lexer}} $\longrightarrow$
\colorbox{amarillo}{\textbf{Parser}} $\longrightarrow$
\colorbox{amarillo}{\textbf{Sem\'antica}} $\longrightarrow$
\colorbox{amarillo}{\textbf{TAC}} $\longrightarrow$
\colorbox{verde}{\textbf{MIPS32}}
\caption{Pipeline completo de la versi\'on 3.0.0.}
\end{figure}

El principio de bloqueo evita que un error temprano se convierta en muchos errores derivados: lexer y parser contin\'uan recolectando diagn\'osticos, pero TAC y MIPS no se ejecutan si el programa no posee significado v\'alido.

\section{C\'odigo de tres direcciones}

TAC descompone expresiones complejas en instrucciones con a lo sumo dos operandos fuente y un resultado. Por ejemplo, \codigo{x = a + b * 4} se convierte en:

\begin{verbatim}
t0 = b * 4
t1 = a + t0
x = t1
\end{verbatim}

Esta forma simplifica la selecci\'on de instrucciones: cada cu\'adruplo tiene una traducci\'on local y los metadatos de tipo determinan si debe usarse la ALU entera o el coprocesador flotante.

\section{Arquitectura MIPS32 y MARS}

MIPS emplea un conjunto reducido de instrucciones y registros expl\'icitos. El proyecto usa registros temporales \codigo{\$t0--\$t9}, argumentos \codigo{\$a0--\$a1}, retornos \codigo{\$v0}/\codigo{\$f0}, stack pointer \codigo{\$sp}, frame pointer \codigo{\$fp}, return address \codigo{\$ra} y el registro \codigo{\$s7} como enlace al entorno de un closure.

MARS 4.5 ensambla MIPS32 y ofrece syscalls educativas. Se emplean syscall 1 para enteros, 2 para flotantes, 4 para strings, 9 para memoria din\'amica y 10 para terminar el programa.

\section{Registros de activaci\'on}

Cada llamada necesita un frame independiente para soportar recursi\'on. El tama\~no proviene de \codigo{ActivationRecord.totalBytes} y se alinea a ocho bytes.

\begin{figure}[ht]
\centering
\renewcommand{\arraystretch}{1.5}
\begin{tabular}{|>{\centering\arraybackslash}p{9cm}|}\hline
\colorbox{amarillo}{\parbox{8.5cm}{\centering\textbf{argumentos del llamador}}}\\\hline
\codigo{\$ra} guardado\\\hline
\codigo{\$fp} anterior\\\hline
\codigo{\$s7} / entorno anterior\\\hline
par\'ametros y variables locales\\\hline
\colorbox{verde!35}{\parbox{8.5cm}{\centering\textbf{temporales TAC $\leftarrow$ \$fp = \$sp}}}\\\hline
\end{tabular}
\caption{Organizaci\'on conceptual del frame.}
\end{figure}

El llamador apila argumentos de izquierda a derecha. La funci\'on reserva el frame, guarda estado y copia par\'ametros a offsets estables. Una \codigo{CALL} consume solo sus propios argumentos, lo cual resuelve correctamente llamadas anidadas.

\chapter{Dise\~no e implementaci\'on}

\section{Contrato del backend}

\archivo{src/mips/generator.ts} expone \codigo{generateMips(tac)}. El resultado usa estados expl\'icitos: \codigo{not-requested}, \codigo{skipped}, \codigo{completed} y \codigo{failed}. Un ensamblador parcial nunca se presenta como listo.

\begin{longtable}{p{3cm}p{10cm}}
\toprule
Estado & Significado\\\midrule
\endhead
not-requested & la fase no fue solicitada\\
skipped & no existe TAC v\'alido por errores anteriores\\
completed & ensamblador generado y validado\\
failed & el backend detect\'o una inconsistencia\\
\bottomrule
\end{longtable}

\section{Resoluci\'on de almacenamiento}

Los operandos globales se convierten en etiquetas del segmento \codigo{.data}. Los operandos de la funci\'on activa usan su offset relativo a \codigo{\$fp}. Las constantes se materializan con \codigo{li}, \codigo{la} o \codigo{l.s}. Los valores persistentes se almacenan despu\'es de cada cu\'adruplo; as\'i no dependen de que un registro temporal sobreviva a saltos o llamadas.

\section{Asignador getReg}

El asignador administra \codigo{\$t0} a \codigo{\$t7} y reinicia su conjunto de uso en cada instrucci\'on TAC. Esta decisi\'on corresponde a la naturaleza de tres direcciones: los valores vivos entre instrucciones est\'an en slots, mientras que los registros sirven para cargar operandos y producir un resultado local.

\section{Selecci\'on de instrucciones}

\begin{longtable}{p{3cm}p{4.5cm}p{5.5cm}}
\toprule
Familia TAC & MIPS & Observaci\'on\\\midrule
\endhead
MOV & li/la/lw y sw & depende de la clase del operando\\
ADD/SUB/MUL & addu/subu/mul & enteros de 32 bits\\
DIV/MOD & div con mflo/mfhi & cociente o residuo\\
EQ--GE & seq, sne, slt, sle, sgt, sge & resultado booleano 0/1\\
IF/GOTO & beq, bne, j & delay slot expl\'icito\\
CALL & jal o jalr & nombre directo o puntero\\
RETURN & \$v0/\$f0 & salta a un ep\'ilogo com\'un\\
PRINT & syscall & seleccionado por tipo\\
\bottomrule
\end{longtable}

\section{Datos compuestos}

Los arreglos guardan longitud en el primer word y elementos desde el segundo. Los objetos reservan \codigo{ClassLayout.instanceSize}; los campos usan offsets calculados en TAC. Los m\'etodos se convierten en direcciones de etiqueta. Un closure reserva un puntero a funci\'on seguido de sus capturas y usa \codigo{\$s7} como entorno activo.

\section{Strings y flotantes}

Las cadenas constantes se internan para evitar duplicados. La concatenaci\'on solicita memoria y copia bytes; la igualdad compara byte a byte. Los flotantes usan instrucciones \codigo{l.s}, \codigo{s.s}, \codigo{add.s}, \codigo{sub.s}, \codigo{mul.s} y \codigo{div.s}. La promoci\'on entera se realiza con \codigo{mtc1} y \codigo{cvt.s.w}.

\section{Mirilla}

La optimizaci\'on local elimina movimientos de un registro hacia s\'i mismo, ajustes nulos de pila, saltos hacia la etiqueta siguiente y \codigo{nop} consecutivos. Se conserva la salida previa para auditor\'ia y se informa cu\'antas instrucciones se removieron.

\chapter{Interfaz y experiencia de uso}

El flujo normal para Entrada v\'alida y Archivo propio solicita modo \codigo{mips}; por tanto, un programa correcto llega autom\'aticamente al backend. El modo TAC permanece disponible como demostraci\'on aislada.

El pipeline visual usa el estado real de cada fase. MIPS aparece como omitido ante errores previos, fallido ante diagn\'osticos del backend y listo para MARS solo en \codigo{completed}. El panel muestra estado, instrucciones finales, optimizaciones, diagn\'osticos, cantidad TAC y pico de registros.

El selector de archivos est\'a dentro de la GUI y acepta \archivo{.cps}. Los resultados se muestran en la misma interfaz; la consola no es necesaria para la evaluaci\'on funcional.

\chapter{CLI y ejecuci\'on externa}

La CLI comparte \codigo{analyzeInput}; no posee un compilador alterno. El modo MIPS acepta \codigo{--out}:

\begin{verbatim}
npm run cli -- examples/mips/07_recursion.cps \
  --mode mips --out factorial.asm
java -jar Mars4_5.jar nc ae1 factorial.asm
\end{verbatim}

El script \archivo{scripts/run-mars.mjs} busca el JAR en el segundo argumento, en \codigo{MARS\_JAR} o en \archivo{tools/Mars4\_5.jar}. MARS sigue siendo una herramienta externa: el IDE genera y exporta, pero no simula internamente.

\chapter{Pruebas y resultados}

\section{Suite automatizada}

La suite final contiene 176 pruebas en 13 archivos. Ocho corresponden directamente al backend y verifican generaci\'on completa, convenci\'on de pila, arreglos, runtime de strings, llamadas anidadas, closures, bloqueo por errores y mirilla. Las pruebas heredadas contin\'uan verificando lexer, parser, sem\'antica, s\'imbolos, TAC, CFG, UI y la exportaci\'on del nuevo modo MIPS.

\section{Validaci\'on real en MARS}

Se ejecutaron quince programas positivos en MARS 4.5. Algunos resultados representativos son:

\begin{longtable}{p{7cm}p{5cm}}
\toprule
Caso & Salida\\\midrule
aritm\'etica & 14\\
funciones y pila & 42\\
arreglos y foreach & 6\\
strings y booleanos & Hola MIPS; true\\
flotantes & 4.0\\
factorial recursivo & 120\\
closure & 1\\
clase, m\'etodo y llamada anidada & 2; 6; 12\\
\bottomrule
\end{longtable}

El caso de errores mezclados produce diagn\'osticos, deja TAC/MIPS en \codigo{skipped} y no crea un archivo ensamblador.

\section{Comandos de verificaci\'on}

\begin{verbatim}
npm ci
npm run check
npm test
npm run build
npm run cli:mips
npm audit --omit=dev
\end{verbatim}

TypeScript, pruebas y build terminan correctamente; la auditor\'ia de dependencias de producci\'on reporta cero vulnerabilidades.

\chapter{Trazabilidad con el enunciado}

\begin{longtable}{p{6cm}p{2.4cm}p{5cm}}
\toprule
Requisito & Estado & Evidencia\\\midrule
\endhead
ANTLR obligatorio & Cumple & gram\'atica y c\'odigo generado\\
MIPS desde TAC & Cumple & API \codigo{generateMips(tac)}\\
GUI y selector de archivo & Cumple & React/Monaco y men\'u Archivo\\
Resultados dentro de GUI & Cumple & paneles de fases, TAC y MIPS\\
M\'ultiples errores y recuperaci\'on & Cumple & listeners, estrategia ANTLR y visitors\\
No TAC/MIPS ante errores & Cumple & estados y pruebas negativas\\
Mirilla & Cumple & optimizador y m\'etricas\\
Ensamblador ejecutable & Cumple & quince ejecuciones MARS\\
No c\'odigo objeto ni ejecuci\'on interna & Cumple & exportaci\'on a herramienta externa\\
\bottomrule
\end{longtable}

\chapter{Limitaciones y trabajo futuro}

Compiscript y el TAC actual no definen una instrucci\'on \codigo{throw} ni un ABI de excepciones; por ello los marcadores \codigo{TRY\_BEGIN}/\codigo{CATCH\_BEGIN} preservan estructura, pero no realizan unwinding de pila. Tampoco existe una fase global de optimizaci\'on, de acuerdo con el alcance: solo se aplica mirilla local.

Como evoluci\'on podr\'ian incorporarse bounds checks con handlers, asignaci\'on global de registros por liveness y un runner de integraci\'on continua que descargue MARS bajo aceptaci\'on de su licencia. Estas mejoras no son necesarias para el alcance actual.

\chapter{Conclusiones}

El Proyecto 3 integra de forma coherente las fases anteriores y produce ensamblador MIPS32 real desde TAC. La convenci\'on de llamadas permite funciones, anidamiento y recursi\'on; el modelo de heap cubre arreglos, objetos, cadenas y closures; los estados impiden exportar resultados inv\'alidos; y la interfaz hace observable todo el proceso.

La evidencia combina pruebas unitarias, integraci\'on, build de producci\'on y ejecuci\'on en MARS 4.5. Por ello, la afirmaci\'on de compatibilidad no depende solamente de documentaci\'on o de ejemplos de ensamblador, sino de resultados reproducibles.

\appendix
\chapter{Mapa de archivos}

\begin{longtable}{p{6cm}p{7cm}}
\toprule
Archivo & Responsabilidad\\\midrule
\endhead
\archivo{src/mips/generator.ts} & backend, runtime, mirilla y diagn\'osticos\\
\archivo{src/tac/types.ts} & contrato de TAC, frames y layouts\\
\archivo{src/lib/analyze.ts} & orquestaci\'on y bloqueo entre fases\\
\archivo{src/ui/App.tsx} & flujo completo para entrada propia\\
\archivo{src/ui/components/RightDock.tsx} & inspecci\'on y descarga MIPS\\
\archivo{src/cli/run.ts} & modo MIPS y escritura \archivo{.asm}\\
\archivo{scripts/run-mars.mjs} & invocaci\'on externa de MARS\\
\archivo{src/__tests__/mips.test.ts} & regresiones del backend\\
\archivo{examples/mips} & demostraciones y caso de bloqueo\\
\bottomrule
\end{longtable}

\end{document}
